%% notes-tex/031-inhibitor-tolerance-data/sections/4-structure.tex

\section{What Hillenmeyer shares with Vanacloig\secstatus{todo}}
\label{sec:structure}

Both matrices are oriented so that negative is a defect, Hillenmeyer's doses and
generation counts are averaged within a compound, and every Vanacloig condition is
correlated with every partner condition over the shared genes
(\file{cross\_dataset\_similarity.py}).
Three statistics ask three different questions. The cross Spearman of two conditions
asks whether the same genes are sensitive to both. The per-gene mean response across
all conditions asks whether the datasets agree on which genes are generally sick. The
Mantel statistic asks whether they agree on which genes move together: the two
gene-by-gene correlation matrices, built from each dataset's condition profiles, are
compared entry by entry against a null that shuffles gene identity.

%% SOURCE: experiments/031-env-chemgen-inhibitor-tolerance/scripts/cross_dataset_similarity.py ->
%% results/structure_{hom,het}.csv, rendered by notes_tex_tables.py.
\begin{table}[htbp]
\centering
\footnotesize
\caption{Shared response structure between Vanacloig and each Hillenmeyer arm. The
cross Spearman columns summarize the full condition-by-condition matrix; the Mantel
column is the Spearman between the two gene-gene profile-correlation matrices on
1{,}500 sampled genes, with the mean and SD of 20 gene-permutation nulls.}
\label{tab:structure}
%% GENERATED by experiments/031-env-chemgen-inhibitor-tolerance/scripts/notes_tex_tables.py
%% SOURCE: results/structure_{hom,het}.csv from cross_dataset_similarity.py. Do not edit by hand.
\begin{tabular}{lrrrrrrl}
\toprule
arm & conditions & shared genes & cross $\rho$ median & 95th pct & max & mean-response $\rho$ & Mantel $\rho$ (null) \\
\midrule
HOM & 140 & 3{,}550 & 0.010 & 0.075 & 0.222 & 0.118 & 0.068 (-0.003 $\pm$ 0.006) \\
HET & 308 & 3{,}577 & 0.003 & 0.052 & 0.183 & 0.036 & 0.024 (-0.001 $\pm$ 0.007) \\
\bottomrule
\end{tabular}

\end{table}

The shared structure is real and small (Table~\ref{tab:structure}). For HOM the Mantel
statistic is eleven null standard deviations from zero and the per-gene mean response
correlates at 0.118, while the median cross Spearman of a condition pair is 0.010 and
the largest is 0.222. HOM carries about three times the shared structure that HET does
on every statistic, which is the expected direction: HOM is the same perturbation class
as Vanacloig, HET is a dosage halving.

%% SOURCE: experiments/031-env-chemgen-inhibitor-tolerance/scripts/cross_dataset_similarity.py ->
%% cross_similarity_heatmap_hom.svg, from results/cross_spearman_hom.csv.
\tcfig{figures/cross_similarity_heatmap_hom.pdf}{\textbf{(a)} Cross Spearman between
each Vanacloig condition (rows) and the 40 HOM conditions with the largest absolute
correlation to any row (columns), over the 3{,}550 shared genes. The solvent-like
inhibitors (ethanol, isobutanol, furfural, gamma-valerolactone) and myclobutanil share
their best match, basifungin; the phenolic aldehydes and ketones match methyl
methanesulfonate and mechlorethamine.}{fig:heatmap}

\notesec{The exact-match compounds}
Two compounds are shared with each arm, and Table~\ref{tab:shared} gives each pair the
same three tests. The cross Spearman is the correlation of the two per-gene profiles.
The rank places the true match among all partner conditions: if the same compound in
the other dataset is the best-correlated condition, the rank is 1. The hit overlap is
the number of genes in the bottom 5\% of both profiles, against the 8.6 expected by
chance, with a one-sided Fisher test.

%% SOURCE: experiments/031-env-chemgen-inhibitor-tolerance/scripts/cross_dataset_similarity.py ->
%% results/shared_compounds_{hom,het}.csv, doses from dataset_axes_comparison.py ->
%% results/shared_compounds.csv, rendered by notes_tex_tables.py.
\begin{table}[htbp]
\centering
\footnotesize
\caption{Transfer of every compound both datasets dose, plus the sodium acetate versus
acetic acid near-pair, which the graph does not join and which serves as the negative
control. Rank is the position of the true match among all partner conditions by cross
Spearman; hits is the overlap of bottom-5\% genes with the chance expectation in
parentheses.}
\label{tab:shared}
%% GENERATED by experiments/031-env-chemgen-inhibitor-tolerance/scripts/notes_tex_tables.py
%% SOURCE: results/shared_compounds_{hom,het}.csv and shared_compounds.csv from cross_dataset_similarity.py and dataset_axes_comparison.py. Do not edit by hand.
\begin{tabular}{llp{3.6cm}rrrrr}
\toprule
compound (V / H) & arm & dose V / H & genes & $\rho$ & rank & hits (exp.) & Fisher $p$ \\
\midrule
benomyl & HOM & 10.0 $\mu$g/mL / 10.0 $\mu$g/mL & 3{,}428 & -0.007 & 88 of 140 & 16 (8.6) & 1.1e-02 \\
methyl methanesulfonate & HOM & fixed / 0.002 \% v/v; 0.004 \% v/v & 3{,}426 & 0.132 & 3 of 140 & 39 (8.6) & 8.8e-17 \\
sodium acetate / acetic acid & HOM & IC30 / IC30 vs acid & 3{,}256 & -0.027 & 121 of 140 & 4 (8.2) & 9.7e-01 \\
benomyl & HET & 10.0 $\mu$g/mL / 13.8 $\mu$M; 2.4 $\mu$M; 27.6 $\mu$M; 3.4 $\mu$M; 6.9 $\mu$M & 3{,}436 & 0.024 & 30 of 308 & 5 (8.6) & 9.4e-01 \\
ferulic acid & HET & IC30 / 3.95 $\mu$M & 3{,}416 & 0.005 & 110 of 308 & 8 (8.6) & 6.3e-01 \\
\bottomrule
\end{tabular}

\end{table}

Figure~\ref{fig:scatter} shows what the table summarizes. Each panel is one shared
compound; each point is one gene, placed at its Vanacloig response (x) and its HOM
response (y), both oriented so that lower is sicker. A gene that the compound hurts in
both screens falls toward the lower left, and the red points are the genes in the
bottom 5\% of both. Methyl methanesulfonate (panel d) is what transfer looks like: a
cloud that leans toward the lower left, 39 genes hit in both against 8.6 expected, and
the true match third of 140 HOM conditions, behind two other DNA-damaging agents.
Benomyl (panel c) is what a hit-level-only transfer looks like: the cloud is round, the
correlation is zero, and the 16 doubly hit genes are the only structure. Sodium acetate
against acetic acid (panel e) is what no transfer looks like: 4 shared hits, fewer than
chance.

%% SOURCE: experiments/031-env-chemgen-inhibitor-tolerance/scripts/cross_dataset_similarity.py ->
%% shared_compound_scatter_hom.svg, from the served responses of the two datasets.
\tcfig{figures/shared_compound_scatter_hom.pdf}{Per-gene responses for the three HOM
pairs, one point per shared gene, bottom-5\% hits in both datasets in red.
\textbf{(c)} Benomyl at 10\,$\mu$g/mL in both. \textbf{(d)} Methyl methanesulfonate,
fixed dose in Vanacloig against 0.002 and 0.004\% in HOM. \textbf{(e)} Sodium acetate
at IC30 in anaerobic SynBase against acetic acid in YPD.}{fig:scatter}

\notesec{The best matches group by chemistry without a shared compound}
The full ranking (\file{results/top\_matches\_hom.csv}) puts basifungin first for
ethanol, isobutanol, furfural, gamma-valerolactone, sodium glyoxylate and myclobutanil,
with myriocin second for furfural and glyoxylate; acetovanillone, acetosyringone and
the phenolic amides match methyl methanesulfonate and mechlorethamine first. Hypothesis
(untested): the solvent-like inhibitors share a membrane or sphingolipid stress profile
and the aldehyde and ketone phenolics share a DNA-damage profile. Section~\ref{sec:chemistry}
asks the general form of this question with molecular encoders.
